\documentclass[10pt,a4paper]{article}
\usepackage[utf8]{inputenc}
\usepackage[T1]{fontenc}
\usepackage[english,spanish]{babel}
\usepackage{geometry}
\geometry{margin=1.8cm}
\usepackage{enumitem}
\usepackage{booktabs}
\usepackage{fancyhdr}
\usepackage{lastpage}
\usepackage{tcolorbox}
\usepackage{listings}
\usepackage{xcolor}

\lstset{
  language=Python,
  basicstyle=\ttfamily\small,
  keywordstyle=\color{blue}\bfseries,
  stringstyle=\color{red},
  commentstyle=\color{gray}\itshape,
  numbers=left,
  numberstyle=\tiny\color{gray},
  numbersep=5pt,
  frame=single,
  breaklines=true,
  showstringspaces=false,
  tabsize=4
}

\pagestyle{fancy}
\fancyhf{}
\rfoot{\thepage/\pageref{LastPage}}
\lhead{\small Python: De Cero a Clases en 7 D\'ias}
\rhead{\small MPrometeo}
\renewcommand{\headrulewidth}{0.4pt}
\setlength{\parindent}{0pt}
\setlength{\parskip}{5pt}

\begin{document}

\begin{center}
{\LARGE \textbf{Python: De Cero a Dominio de Clases}}\\[4pt]
{\large Gu\'ia ultra-compacta en 7 d\'ias}\\[2pt]
\rule{\textwidth}{0.8pt}
\end{center}

% ============================================================
\section*{D\'ia 1: Variables, Tipos, Operadores}
% ============================================================

\begin{lstlisting}
# Variables: no necesitas declarar tipo
nombre = "MPrometeo"    # str
edad = 25               # int
pi = 3.14159            # float
es_estudiante = True    # bool

# Operadores basicos
suma = 10 + 3           # 13
potencia = 2 ** 10      # 1024
modulo = 17 % 5         # 2
division_entera = 17 // 5  # 3

# Strings
saludo = f"Hola {nombre}, tienes {edad} anios"
print(saludo)
print(len(nombre))      # 9
print(nombre.upper())   # MPROMETEO
print(nombre[0:4])      # MPro (slicing)

# Input
x = input("Dame un numero: ")
x = int(x)  # convertir str a int

# Type checking
print(type(pi))         # <class 'float'>
\end{lstlisting}

\textbf{Memorizar:} \texttt{int}, \texttt{float}, \texttt{str},
\texttt{bool}, \texttt{f-strings}, \texttt{type()}, slicing
\texttt{[inicio:fin]}.

% ============================================================
\section*{D\'ia 2: Condicionales y Bucles}
% ============================================================

\begin{lstlisting}
# if / elif / else
nota = 85
if nota >= 90:
    print("A")
elif nota >= 80:
    print("B")
elif nota >= 70:
    print("C")
else:
    print("F")

# Operadores logicos: and, or, not
if edad >= 18 and es_estudiante:
    print("Adulto estudiante")

# Ternario (una linea)
resultado = "Aprobado" if nota >= 60 else "Reprobado"

# for loop
for i in range(5):       # 0,1,2,3,4
    print(i)

for letra in "Pynchon":  # itera sobre string
    print(letra)

# while loop
contador = 0
while contador < 5:
    print(contador)
    contador += 1

# break y continue
for i in range(10):
    if i == 3:
        continue  # salta el 3
    if i == 7:
        break     # para en 7
    print(i)      # imprime 0,1,2,4,5,6

# enumerate (indice + valor)
nombres = ["Alice", "Bob", "Charlie"]
for i, nombre in enumerate(nombres):
    print(f"{i}: {nombre}")
\end{lstlisting}

\textbf{Memorizar:} indentaci\'on = 4 espacios (obligatorio),
\texttt{range(inicio, fin, paso)}, \texttt{enumerate()}.

% ============================================================
\section*{D\'ia 3: Listas, Tuplas, Diccionarios, Sets}
% ============================================================

\begin{lstlisting}
# LISTAS (mutables, ordenadas)
frutas = ["manzana", "banana", "cereza"]
frutas.append("durazno")     # agregar al final
frutas.insert(1, "kiwi")     # insertar en posicion
frutas.remove("banana")      # eliminar por valor
ultimo = frutas.pop()        # eliminar y devolver ultimo
frutas.sort()                # ordenar in-place
print(frutas[0])             # primer elemento
print(frutas[-1])            # ultimo elemento
print(frutas[1:3])           # slicing

# List comprehension (IMPORTANTISIMO)
cuadrados = [x**2 for x in range(10)]
pares = [x for x in range(20) if x % 2 == 0]

# TUPLAS (inmutables, ordenadas)
coordenada = (10, 20)
x, y = coordenada  # desempaquetado

# DICCIONARIOS (key:value, mutables)
alumno = {
    "nombre": "MPrometeo",
    "edad": 25,
    "carrera": "Ciencias"
}
print(alumno["nombre"])
alumno["gpa"] = 9.5           # agregar key
del alumno["edad"]             # eliminar key
print(alumno.get("email", "N/A"))  # default si no existe

# Iterar diccionario
for key, value in alumno.items():
    print(f"{key}: {value}")

# SETS (unicos, sin orden)
a = {1, 2, 3, 4}
b = {3, 4, 5, 6}
print(a | b)   # union: {1,2,3,4,5,6}
print(a & b)   # interseccion: {3,4}
print(a - b)   # diferencia: {1,2}
\end{lstlisting}

\textbf{Memorizar:} List comprehension es ESENCIAL. Diccionarios son
la estructura m\'as usada en Python real.

% ============================================================
\section*{D\'ia 4: Funciones}
% ============================================================

\begin{lstlisting}
# Funcion basica
def saludar(nombre):
    return f"Hola, {nombre}!"

print(saludar("MPrometeo"))

# Parametros por defecto
def potencia(base, exp=2):
    return base ** exp

print(potencia(3))      # 9
print(potencia(3, 3))   # 27

# *args (argumentos variables)
def sumar(*numeros):
    return sum(numeros)

print(sumar(1, 2, 3, 4))  # 10

# **kwargs (keyword arguments variables)
def perfil(**datos):
    for k, v in datos.items():
        print(f"{k}: {v}")

perfil(nombre="MPrometeo", edad=25, lenguaje="Python")

# Lambda (funcion anonima en una linea)
doble = lambda x: x * 2
print(doble(5))  # 10

# map, filter
nums = [1, 2, 3, 4, 5]
dobles = list(map(lambda x: x*2, nums))
pares = list(filter(lambda x: x%2==0, nums))

# Funcion que devuelve funcion (closure)
def multiplicador(n):
    def multiplicar(x):
        return x * n
    return multiplicar

triple = multiplicador(3)
print(triple(10))  # 30
\end{lstlisting}

\textbf{Memorizar:} \texttt{*args}, \texttt{**kwargs}, \texttt{lambda},
\texttt{map()}, \texttt{filter()}, closures.

% ============================================================
\section*{D\'ia 5: Archivos, Excepciones, M\'odulos}
% ============================================================

\begin{lstlisting}
# Leer archivo
with open("datos.txt", "r") as f:
    contenido = f.read()
    # o linea por linea:
    # for linea in f:
    #     print(linea.strip())

# Escribir archivo
with open("salida.txt", "w") as f:
    f.write("Hola mundo\n")

# Excepciones
try:
    resultado = 10 / 0
except ZeroDivisionError:
    print("No se puede dividir entre cero")
except Exception as e:
    print(f"Error: {e}")
finally:
    print("Esto siempre se ejecuta")

# Levantar excepciones propias
def dividir(a, b):
    if b == 0:
        raise ValueError("Divisor no puede ser cero")
    return a / b

# Modulos
import math
from random import randint, choice
import os

print(math.sqrt(144))     # 12.0
print(randint(1, 100))    # numero aleatorio
print(os.getcwd())        # directorio actual

# Crear tu propio modulo: guarda funciones en un .py
# e importalo con: import mi_modulo
\end{lstlisting}

\textbf{Memorizar:} \texttt{with open()} siempre (cierra archivo
autom\'aticamente), \texttt{try/except/finally}.

% ============================================================
\section*{D\'ia 6: Programaci\'on Orientada a Objetos (POO)}
% ============================================================

\begin{lstlisting}
# Clase basica
class Estudiante:
    # Constructor
    def __init__(self, nombre, edad, carrera):
        self.nombre = nombre    # atributo de instancia
        self.edad = edad
        self.carrera = carrera
        self.materias = []      # lista vacia por defecto

    # Metodo
    def inscribir(self, materia):
        self.materias.append(materia)
        print(f"{self.nombre} inscrito en {materia}")

    # Metodo con return
    def es_adulto(self):
        return self.edad >= 18

    # Representacion string
    def __str__(self):
        return f"{self.nombre} ({self.carrera})"

    # Representacion para debugging
    def __repr__(self):
        return f"Estudiante('{self.nombre}', {self.edad})"

# Crear objetos (instancias)
alumno1 = Estudiante("MPrometeo", 25, "Ciencias")
alumno2 = Estudiante("Alice", 20, "Fisica")

alumno1.inscribir("TOEFL")
alumno1.inscribir("Python")
print(alumno1)               # MPrometeo (Ciencias)
print(alumno1.materias)      # ['TOEFL', 'Python']
print(alumno1.es_adulto())   # True
\end{lstlisting}

\subsection*{Herencia}

\begin{lstlisting}
class Persona:
    def __init__(self, nombre, edad):
        self.nombre = nombre
        self.edad = edad

    def presentarse(self):
        return f"Soy {self.nombre}, tengo {self.edad}"

# Herencia: Estudiante ES UNA Persona
class Estudiante(Persona):
    def __init__(self, nombre, edad, carrera):
        super().__init__(nombre, edad)  # llama al padre
        self.carrera = carrera

    # Override (sobreescribir metodo del padre)
    def presentarse(self):
        base = super().presentarse()
        return f"{base}, estudio {self.carrera}"

class Profesor(Persona):
    def __init__(self, nombre, edad, departamento):
        super().__init__(nombre, edad)
        self.departamento = departamento

est = Estudiante("MPrometeo", 25, "Ciencias")
prof = Profesor("Aldana", 55, "Complejidad")
print(est.presentarse())
# Soy MPrometeo, tengo 25, estudio Ciencias
print(prof.presentarse())
# Soy Aldana, tengo 55
\end{lstlisting}

\subsection*{Encapsulamiento}

\begin{lstlisting}
class CuentaBancaria:
    def __init__(self, titular, saldo=0):
        self.titular = titular
        self.__saldo = saldo  # __ = privado (name mangling)

    def depositar(self, cantidad):
        if cantidad > 0:
            self.__saldo += cantidad

    def retirar(self, cantidad):
        if 0 < cantidad <= self.__saldo:
            self.__saldo -= cantidad
        else:
            print("Fondos insuficientes")

    @property
    def saldo(self):  # getter como propiedad
        return self.__saldo

cuenta = CuentaBancaria("MPrometeo", 1000)
cuenta.depositar(500)
print(cuenta.saldo)  # 1500 (usa @property)
# cuenta.__saldo     # ERROR: atributo privado
\end{lstlisting}

% ============================================================
\section*{D\'ia 7: Clases Avanzadas + Proyecto Final}
% ============================================================

\subsection*{M\'etodos de clase y est\'aticos}

\begin{lstlisting}
class Matematicas:
    pi = 3.14159  # atributo de clase (compartido)

    @staticmethod  # no necesita self ni cls
    def sumar(a, b):
        return a + b

    @classmethod   # recibe la clase, no la instancia
    def area_circulo(cls, radio):
        return cls.pi * radio ** 2

print(Matematicas.sumar(3, 5))        # 8
print(Matematicas.area_circulo(10))   # 314.159
\end{lstlisting}

\subsection*{Dunder Methods (m\'etodos m\'agicos)}

\begin{lstlisting}
class Vector:
    def __init__(self, x, y):
        self.x = x
        self.y = y

    def __add__(self, otro):       # v1 + v2
        return Vector(self.x + otro.x, self.y + otro.y)

    def __mul__(self, escalar):    # v1 * 3
        return Vector(self.x * escalar, self.y * escalar)

    def __len__(self):             # len(v1)
        return int((self.x**2 + self.y**2)**0.5)

    def __eq__(self, otro):        # v1 == v2
        return self.x == otro.x and self.y == otro.y

    def __str__(self):
        return f"({self.x}, {self.y})"

v1 = Vector(3, 4)
v2 = Vector(1, 2)
v3 = v1 + v2
print(v3)       # (4, 6)
print(len(v1))  # 5
print(v1 * 2)   # (6, 8)
\end{lstlisting}

\subsection*{Proyecto Final: Sistema de Calificaciones TOEFL}

\begin{lstlisting}
class TOEFLScore:
    # Tabla de conversion (simplificada)
    CEFR_MAP = {
        "A2": (337, 459),
        "B1": (460, 542),
        "B2": (543, 626),
        "C1": (627, 677)
    }

    def __init__(self, nombre, listening, structure, reading):
        self.nombre = nombre
        self.listening = listening  # scaled 31-68
        self.structure = structure
        self.reading = reading

    @property
    def total(self):
        return round((self.listening + self.structure
                       + self.reading) / 3 * 10)

    @property
    def cefr(self):
        for nivel, (bajo, alto) in self.CEFR_MAP.items():
            if bajo <= self.total <= alto:
                return nivel
        if self.total < 337:
            return "A1"
        return "C1+"

    def __str__(self):
        return (f"{self.nombre}: L={self.listening} "
                f"S={self.structure} R={self.reading} "
                f"| Total={self.total} | CEFR={self.cefr}")

    def __gt__(self, otro):  # comparar scores
        return self.total > otro.total

    @classmethod
    def from_raw(cls, nombre, l_raw, s_raw, r_raw):
        """Crea desde raw scores (simplificado)."""
        # Conversion super simplificada
        l = 31 + int(l_raw / 50 * 37)
        s = 31 + int(s_raw / 40 * 37)
        r = 31 + int(r_raw / 50 * 36)
        return cls(nombre, l, s, r)

# Uso
yo = TOEFLScore("MPrometeo", 49, 52, 54)
print(yo)
# MPrometeo: L=49 S=52 R=54 | Total=517 | CEFR=B1

mock1 = TOEFLScore.from_raw("Mock1", 35, 28, 40)
mock2 = TOEFLScore.from_raw("Mock2", 40, 32, 43)
print(mock1)
print(mock2)
print(f"Mock2 mejor que Mock1? {mock2 > mock1}")

# Lista de scores
scores = [yo, mock1, mock2]
scores.sort(key=lambda s: s.total, reverse=True)
for s in scores:
    print(s)
\end{lstlisting}

% ============================================================
\section*{Resumen de Conceptos por D\'ia}
% ============================================================

\begin{center}
\small
\begin{tabular}{cl}
\toprule
\textbf{D\'ia} & \textbf{Conceptos clave} \\
\midrule
1 & Variables, tipos, f-strings, slicing, type() \\
2 & if/elif/else, for, while, break, continue, enumerate \\
3 & Listas, tuplas, diccionarios, sets, list comprehension \\
4 & Funciones, *args, **kwargs, lambda, map, filter, closures \\
5 & Archivos (with open), try/except, import, modulos \\
6 & class, \_\_init\_\_, self, herencia, super(), @property \\
7 & @staticmethod, @classmethod, dunder methods, proyecto \\
\bottomrule
\end{tabular}
\end{center}

\end{document}